\phantomsection
\section*{Заключение}
\addcontentsline{toc}{section}{Заключение}

В ходе выполнения курсовой работы было спроектировано и реализовано веб-приложение для учёта комплектующих, сборки и продажи персональных компьютеров. Поставленная цель достигнута полностью: создан готовый к внедрению программный продукт, покрывающий бизнес-процессы малой компании по продаже комплектующих и сборке ПК под заказ.

\textbf{В первой главе} проведён анализ предметной области, выполнен сравнительный обзор существующих программных продуктов (DNS, «Регард», «НИКС», «Ситилинк») и универсальных CRM/ERP-систем. Установлено, что распространённые системы ориентированы преимущественно на витрину товаров и не покрывают в комплексе задачи внутреннего учёта комплектующих, конфигурирования сборок и аналитики для малого бизнеса. Анализ программных инструментов разработки позволил обоснованно выбрать технологический стек на основе Python, Django, Bootstrap и PostgreSQL.

\textbf{Во второй главе} выполнено проектирование разрабатываемого веб-приложения: проведён анализ целевой аудитории и трёх пользовательских ролей (администратор, менеджер, клиент), построена матрица прав доступа CRUD, описаны варианты использования и оформлены 16 пользовательских историй, спроектирована модель данных на восемь сущностей с описанием связей и ограничений целостности, разработаны wireframe-макеты основных страниц приложения.

\textbf{В третьей главе} выполнена программная реализация. Разработана модульная структура Django-проекта \texttt{pcshop} из четырёх приложений, реализованы все шесть задач раздела~3:
\begin{enumerate}
    \item базовая структура приложения и шаблоны страниц на Bootstrap~5 с современным дизайн-кодом (Inter, градиентные элементы, sticky-навбар, palette \texttt{\#4F46E5}/\texttt{\#F59E0B});
    \item аутентификация по email и ролевая модель доступа с миксинами и декораторами;
    \item полноценные CRUD-интерфейсы для каталога, сборок, заказов и пользователей с пагинацией, фильтрами, поиском и inline-formset для сборок;
    \item инновационный интерактивный конфигуратор сборки ПК на Alpine.js с пошаговым мастером из восьми шагов, проверкой совместимости комплектующих (4 правила: сокеты CPU--MB, тип памяти MB--RAM, мощность БП, TDP кулера), расчётом потребляемой мощности и сохранением состояния в \texttt{localStorage};
    \item аналитический модуль с агрегированием данных через выражения \texttt{Sum} и \texttt{F}: выручка, прибыль, средний чек, топ-10 товаров и контроль остатков с цветовой индикацией;
    \item модуль импорта и экспорта файлов в формате CSV для каталога и отчётности по заказам с поддержкой BOM для корректной работы с Microsoft Excel.
\end{enumerate}

\textbf{Практическая значимость} работы заключается в создании готового к внедрению программного продукта, который может быть использован действующими магазинами комплектующих и небольшими сборщиками ПК для автоматизации ключевых бизнес-процессов и снижения операционных издержек. Архитектура приложения позволяет масштабироваться по мере роста бизнеса (поддерживается переключение на PostgreSQL без изменения кода через переменные окружения).

\textbf{Особенностью реализации} является конфигуратор сборки ПК, который в отличие от рассмотренных в первой главе аналогов (DNS, НИКС, «Регард») сочетает в себе:
\begin{itemize}
    \item реактивную проверку совместимости в момент выбора каждого компонента (без перезагрузки страницы);
    \item автоматическую рекомендацию оптимальной по цене конфигурации;
    \item расчёт энергопотребления с цветовой индикацией запаса по мощности;
    \item возможность одновременно положить готовую сборку в корзину как единый товарный позицию;
    \item сохранение прогресса в \texttt{localStorage} браузера для возврата к сборке в любой момент.
\end{itemize}

\textbf{Перспективы развития.} Дальнейшее развитие проекта предполагает:
\begin{itemize}
    \item расширение правил совместимости комплектующих (поддержка стандартов PCI Express, охлаждения корпуса, форм-факторов материнских плат);
    \item интеграцию платёжных систем (СБП, банковский эквайринг) и системы email-уведомлений на этапах оформления и смены статуса заказа;
    \item подключение картографических сервисов для расчёта стоимости и сроков доставки;
    \item внедрение системы предиктивной аналитики для прогнозирования остатков и спроса;
    \item развёртывание приложения в облачной инфраструктуре с использованием Docker и Nginx.
\end{itemize}

При разработке пояснительной записки к курсовой работе строго соблюдались требования~\cite{gost_report}, определяющие структуру и правила оформления научно-исследовательских работ.

Исходный код проекта размещён в открытом репозитории GitHub:
\begin{center}
    \url{https://github.com/Haskeri/kpweb}
\end{center}

Развёрнутое веб-приложение доступно по адресу:
\begin{center}
    \url{https://exam.anufriev.pvpcult.ru/}
\end{center}
